\hfuzz 100pt % lässt die overfull hbox warnings verschwinden
\newcommand{\noa}[1]{~{\color{red}Noa: #1}}
\newcommand{\nick}[1]{~{\color{blue}Nick: #1}}

%%
%% Packages
%%

\usepackage{xltabular} % bessere Tabellen
\usepackage{comment} % für Kommentare
\usepackage{scalerel} % richtige Skalierung der speziellen Kategorienzeichen
\RequirePackage{enumerate}
\RequirePackage{mathrsfs}
\RequirePackage{euscript}
\usepackage[shortlabels]{enumitem}
\usepackage{array,longtable,booktabs,cancel,ragged2e,multirow}
\usepackage{titlesec} % Titelanpassungen
\RequirePackage[autostyle=true]{csquotes}
\usepackage{graphicx,wrapfig}
\usepackage{tikz}
\usepackage{tikz-cd}
\usepackage{quiver} % kommutative Diagramme

%%
%% Categories
%%

\newcommand{\fin}{\mathrm{Fin}}
\let\Fin\fin
\newcommand{\finert}{{\operatorname{Fin}_{\ast,\mathrm{inert}}}}
\newcommand{\inert}{{\bbGamma_{\ast,\mathrm{inert}}}}
\newcommand{\set}{\mathrm{Set}}
\let\Set\set
\newcommand{\Ab}{\mathrm{Ab}}
\newcommand{\Grp}{\mathrm{Grp}}
\newcommand{\Ring}{\mathrm{Ring}}
\newcommand{\Mod}{\mathrm{Mod}} % \Mod_R
\newcommand{\ModR}{\Mod_R}
\newcommand{\Vect}{\mathrm{Vect}} % \Vect(\R), \Vect_n(\R)
\newcommand{\Mon}{\mathrm{Mon}}
\newcommand{\cMon}{\mathrm{CMon}}
\newcommand{\bool}{\mathrm{Bool}}
\newcommand{\sbool}{\mathrm{Bool}_\sigma}
\newcommand{\cbool}{\mathrm{Bool}_{\mathrm{comp}}}

\newcommand{\Ani}{\mathrm{Ani}}

\newcommand{\Top}{\mathrm{Top}}
\newcommand{\stone}{\mathrm{Stone}}
\newcommand{\CHaus}{\mathrm{CHaus}}
\newcommand{\extr}{\mathrm{Extr}}
\newcommand{\prof}{\Pro(\Fin)}
\newcommand{\prolak}{\prof_\kappa}
\newcommand{\LCA}{\mathrm{LCA}}

\newcommand{\sSet}{\mathrm{sSet}}
\newcommand{\sAb}{\mathrm{sAb}}

\newcommand{\Cat}{\mathrm{Cat}}
\newcommand{\PrL}{\mathrm{Pr}^\mathrm{L}}

\newcommand{\cond}{\mathrm{Cond}}
\let\Cond\cond
\newcommand{\CondSet}{\mathrm{CondSet}}
\newcommand{\CondAni}{\mathrm{CondAni}}
\newcommand{\cAb}{\mathrm{CondAb}}
\newcommand{\CondAb}{\mathrm{CondAb}}
\newcommand{\CondRing}{\mathrm{CondRing}}
\newcommand{\CondMod}{\mathrm{CondMod}}
\newcommand{\qcqs}{\mathrm{qcqs}}
\newcommand{\liq}{\mathrm{liq}}
\newcommand{\MMp}{\mathcal{M}_{<p}}

\newcommand{\cSet}{\mathrm{cSet}}
\newcommand{\ccSet}{\mathrm{cSet}^\mathrm{conc}}
\newcommand{\cSetnil}{{\mathrm{cSet}_\mathrm{fib}}}
\newcommand{\ccSetnil}{\mathrm{cSet}_\mathrm{fib}^\mathrm{conc}}
\newcommand{\cCond}{\mathrm{cCond}}


%%
%% Operators
%%

\newcommand{\Id}{\operatorname{Id}}
\newcommand{\id}{\operatorname{id}}

\newcommand{\subobj}{\operatorname{Sub}}
\newcommand{\su}{\operatorname{Sub}_{\bbDelta}}

\newcommand{\Gprod}{{\prod^{\mathrm{Gray}}}}
\newcommand{\ev}{\operatorname{ev}}
\newcommand{\Bun}{\operatorname{Bun}}
\newcommand{\HK}{\operatorname{HK}}
\newcommand{\Fun}{\operatorname{Fun}}
\newcommand{\PSh}{\operatorname{PSh}}
\newcommand{\Sh}{\mathrm{Sh}}
\newcommand{\EM}{\operatorname{EM}}
\newcommand{\Ind}{\operatorname{Ind}}
\newcommand{\Pro}{\operatorname{Pro}}
\newcommand{\sInd}{\operatorname{sInd}}
\newcommand{\coker}{\operatorname{coker}}
\newcommand{\Lan}{\operatorname{Lan}}
\newcommand{\Ran}{\operatorname{Ran}}
\newcommand{\Coim}{\operatorname{coim}}
\newcommand{\Ch}{\operatorname{Ch}}
\newcommand{\tot}{\operatorname{tot}}
\newcommand{\co}{\operatorname{co}}
\newcommand{\codim}{\operatorname{codim}}
\newcommand{\dist}{\operatorname{dist}}
\newcommand{\divop}{\operatorname{div}}
\newcommand{\dom}{\operatorname{dom}}
\newcommand{\esssup}{\operatorname{ess\,sup}}
\newcommand{\essrang}{\operatorname{ess\,rang}}
\newcommand{\ext}{\operatorname{ext}}
\newcommand{\fix}{\operatorname{fix}}
\renewcommand{\gcd}{\operatorname{gcd}}
\newcommand{\grad}{\operatorname{grad}}
\newcommand{\interior}{\operatorname{int}}
\newcommand{\lcm}{\operatorname{lcm}}
\newcommand{\lin}{\operatorname{lin}}
\newcommand{\orb}{\operatorname{orb}}
\newcommand{\proj}{\operatorname{proj}}
\newcommand{\ran}{\operatorname{ran}}
\newcommand{\rank}{\operatorname{rank}}
\newcommand{\res}{\operatorname{res}}
\newcommand{\rot}{\operatorname{rot}}
\newcommand{\supp}{\operatorname{supp}}
\newcommand{\tr}{\operatorname{tr}}
\newcommand{\sign}{\operatorname{sign}}
\newcommand{\diag}{\operatorname{diag}}
\DeclareMathOperator*{\argmin}{arg\,min}
\DeclareMathOperator*{\argmax}{arg\,max}
\newcommand{\cosk}{\operatorname{cosk}}
\newcommand{\sing}{\operatorname{Sing}}
\newcommand{\cof}{\operatorname{cof}}

\newcommand{\im}{\operatorname{im}}
\renewcommand{\Im}{\operatorname{Im}}
\renewcommand{\Re}{\operatorname{Re}}
\newcommand{\Var}{\operatorname{Var}}
\newcommand{\Cov}{\operatorname{Cov}}

\newcommand{\Hom}{\operatorname{Hom}}
\let\hom\Hom
\newcommand{\ihom}{\underline{\hom}}
\newcommand{\iHom}{\underline{\hom}}
\newcommand{\Rhom}{\operatorname{RHom}}
\newcommand{\RHom}{\Rhom}
\newcommand{\iRHom}{\underline{\Rhom}}
\newcommand{\Ext}{\operatorname{Ext}}
\newcommand{\iExt}{\underline{\Ext}}
\newcommand{\Iso}{\operatorname{Iso}}
\newcommand{\Aut}{\operatorname{Aut}}
\newcommand{\Trans}{\operatorname{Trans}}
\newcommand{\End}{\operatorname{End}}
\newcommand{\Mat}{\operatorname{Mat}} % \Mat(n, m; \R), \Mat_n(\R)
\newcommand{\Mor}{\operatorname{Mor}}
\newcommand{\Per}{\operatorname{Per}}
\newcommand{\GL}{\operatorname{GL}} % \GL(n, \R)
\newcommand{\SL}{\operatorname{SL}} % \SL(n, \R)
\newcommand{\SO}{\operatorname{SO}}
\newcommand{\Sym}{\operatorname{Sym}}
\newcommand{\Alt}{\operatorname{Alt}}
\newcommand{\Bil}{\operatorname{Bil}}
\newcommand{\Gal}{\operatorname{Gal}} % \Gal(E/F)

\newcommand{\Ob}{\mathrm{Ob}}
\newcommand{\ob}{\Ob}
\newcommand{\op}{\mathrm{op}}
\newcommand{\from}{\leftarrow}

%%
%% Letters
%%

\newcommand\cA{\mathcal{A}}
\newcommand\cB{\mathcal{B}}
\newcommand\cC{\mathcal{C}}
\newcommand\cD{\mathcal{D}}
\newcommand\cE{\mathcal{E}}
\newcommand\cF{\mathcal{F}}
\newcommand\cG{\mathcal{G}}
\newcommand\cH{\mathcal{H}}
\newcommand\cI{\mathcal{I}}
\newcommand\cJ{\mathcal{J}}
\newcommand\cK{\mathcal{K}}
\newcommand\cL{\mathcal{L}}
\newcommand\cM{\mathcal{M}}
\newcommand\cN{\mathcal{N}}
\newcommand\cO{\mathcal{O}}
\newcommand\cP{\mathcal{P}}
\newcommand\cQ{\mathcal{Q}}
\newcommand\cR{\mathcal{R}}
\newcommand\cS{\mathcal{S}}
\newcommand\cT{\mathcal{T}}
\newcommand\cU{\mathcal{U}}
\newcommand\cV{\mathcal{V}}
\newcommand\cW{\mathcal{W}}
\newcommand\cX{\mathcal{X}}
\newcommand\cY{\mathcal{Y}}
\newcommand\cZ{\mathcal{Z}}

\newcommand\mcA{\cA}
\newcommand\mcB{\cB}
\newcommand\mcC{\cC}
\newcommand\mcD{\cD}
\newcommand\mcE{\cE}
\newcommand\mcF{\cF}
\newcommand\mcG{\cG}
\newcommand\mcH{\cH}
\newcommand\mcI{\cI}
\newcommand\mcJ{\cJ}
\newcommand\mcK{\cK}
\newcommand\mcL{\cL}
\newcommand\mcM{\cM}
\newcommand\mcN{\cN}
\newcommand\mcO{\cO}
\newcommand\mcP{\cP}
\newcommand\mcQ{\cQ}
\newcommand\mcR{\cR}
\newcommand\mcS{\cS}
\newcommand\mcT{\cT}
\newcommand\mcU{\cU}
\newcommand\mcV{\cV}
\newcommand\mcW{\cW}
\newcommand\mcX{\cX}
\newcommand\mcY{\cY}
\newcommand\mcZ{\cZ}

\newcommand{\ua}{\mathrm{a}}
\newcommand{\ub}{\mathrm{b}}
\newcommand{\uc}{\mathrm{c}}
\newcommand{\ud}{\mathrm{d}}
\newcommand{\ue}{\mathrm{e}}
\newcommand{\uf}{\mathrm{f}}
\newcommand{\ug}{\mathrm{g}}
\newcommand{\uh}{\mathrm{h}}
\newcommand{\ui}{\mathrm{i}}
\newcommand{\uj}{\mathrm{j}}
\newcommand{\uk}{\mathrm{k}}
\newcommand{\ul}{\mathrm{l}}
\newcommand{\um}{\mathrm{m}}
\newcommand{\un}{\mathrm{n}}
\newcommand{\uo}{\mathrm{o}}

\AtBeginDocument{
  \ifdefined\up
    \let\frup\up
    \renewcommand{\up}{\mathrm{p}}
  \else
    \newcommand{\up}{\mathrm{p}}
  \fi
}
\newcommand{\uq}{\mathrm{q}}
\newcommand{\ur}{\mathrm{r}}
\newcommand{\us}{\mathrm{s}}
\newcommand{\ut}{\mathrm{t}}
\newcommand{\uu}{\mathrm{u}}
\newcommand{\uv}{\mathrm{v}}
\newcommand{\uw}{\mathrm{w}}
\newcommand{\ux}{\mathrm{x}}
\newcommand{\uy}{\mathrm{y}}
\newcommand{\uz}{\mathrm{z}}

\newcommand{\uA}{\mathrm{A}}
\newcommand{\uB}{\mathrm{B}}
\newcommand{\uC}{\mathrm{C}}
\newcommand{\uD}{\mathrm{D}}
\newcommand{\uE}{\mathrm{E}}
\newcommand{\uF}{\mathrm{F}}
\newcommand{\uG}{\mathrm{G}}
\newcommand{\uH}{\mathrm{H}}
\newcommand{\uI}{\mathrm{I}}
\newcommand{\uJ}{\mathrm{J}}
\newcommand{\uK}{\mathrm{K}}
\newcommand{\uL}{\mathrm{L}}
\newcommand{\uM}{\mathrm{M}}
\newcommand{\uN}{\mathrm{N}}
\newcommand{\uO}{\mathrm{O}}
\newcommand{\uP}{\mathrm{P}}
\newcommand{\uQ}{\mathrm{Q}}
\newcommand{\uR}{\mathrm{R}}
\newcommand{\uS}{\mathrm{S}}
\newcommand{\uT}{\mathrm{T}}
\newcommand{\uU}{\mathrm{U}}
\newcommand{\uV}{\mathrm{V}}
\newcommand{\uW}{\mathrm{W}}
\newcommand{\uX}{\mathrm{X}}
\newcommand{\uY}{\mathrm{Y}}
\newcommand{\uZ}{\mathrm{Z}}


\newcommand{\N}{{\mathbb{N}}}
\newcommand{\Z}{\mathbb{Z}}
\newcommand{\Q}{\mathbb{Q}}
\newcommand{\R}{\mathbb{R}}
\newcommand{\C}{\mathbb{C}}
\newcommand{\D}{\mathbb{D}}
\newcommand{\K}{\mathbb{K}}
\newcommand{\T}{\mathbb{T}}
\newcommand{\F}{\mathbb{F}}
\renewcommand{\S}{\mathbb{S}}
\renewcommand{\O}{\mathcal{O}}
\renewcommand{\P}{\mathbb{P}}
\newcommand{\E}{\mathbb{E}}

\newcommand{\bS}{\mathbf{S}}
\newcommand{\bM}{\mathbf{M}}
\newcommand{\fM}{\mathfrak{M}}

%%
%% Measures
%%

\newcommand{\da}{\,\mathrm{d}a}
\newcommand{\db}{\,\mathrm{d}b}
\newcommand{\dc}{\,\mathrm{d}c}
\newcommand{\dd}{\,\mathrm{d}d}
\newcommand{\de}{\,\mathrm{d}e}
\newcommand{\df}{\,\mathrm{d}f}
\newcommand{\dg}{\,\mathrm{d}g}
\newcommand{\di}{\,\mathrm{d}i}
\newcommand{\dk}{\,\mathrm{d}k}
\newcommand{\dl}{\,\mathrm{d}l}
\newcommand{\dm}{\,\mathrm{d}m}
\newcommand{\dn}{\,\mathrm{d}n}
\newcommand{\dr}{\,\mathrm{d}r}
\newcommand{\ds}{\,\mathrm{d}s}
\newcommand{\dt}{\,\mathrm{d}t}
\newcommand{\du}{\,\mathrm{d}u}
\newcommand{\dv}{\,\mathrm{d}v}
\newcommand{\dw}{\,\mathrm{d}w}
\newcommand{\dx}{\,\mathrm{d}x}
\newcommand{\dy}{\,\mathrm{d}y}
\newcommand{\dz}{\,\mathrm{d}z}

\newcommand{\dalpha}{\,\mathrm{d}\alpha}
\newcommand{\dbeta}{\,\mathrm{d}\beta}
\newcommand{\dgamma}{\,\mathrm{d}\gamma}
\newcommand{\ddelta}{\,\mathrm{d}\delta}
\newcommand{\depsilon}{\,\mathrm{d}\epsilon}
\newcommand{\dzeta}{\,\mathrm{d}\zeta}
\newcommand{\deta}{\,\mathrm{d}\eta}
\newcommand{\dtheta}{\,\mathrm{d}\theta}
\newcommand{\diota}{\,\mathrm{d}\iota}
\newcommand{\dkappa}{\,\mathrm{d}\kappa}
\newcommand{\dlambda}{\,\mathrm{d}\lambda}
\newcommand{\dmu}{\,\mathrm{d}\mu}
\newcommand{\dnu}{\,\mathrm{d}\nu}
\newcommand{\dxi}{\,\mathrm{d}\xi}
% No \omikron because it's the same as a plain o.
\newcommand{\dpi}{\,\mathrm{d}\pi}
\newcommand{\drho}{\,\mathrm{d}\rho}
\newcommand{\dsigma}{\,\mathrm{d}\sigma}
\newcommand{\dtau}{\,\mathrm{d}\tau}
\newcommand{\dupsilon}{\,\mathrm{d}\upsilon}
\newcommand{\dphi}{\,\mathrm{d}\phi}
\newcommand{\dchi}{\,\mathrm{d}\chi}
\newcommand{\dpsi}{\,\mathrm{d}\psi}
\newcommand{\domega}{\,\mathrm{d}\omega}

\newcommand{\dA}{\,\mathrm{d}A}
\newcommand{\dB}{\,\mathrm{d}B}
\newcommand{\dC}{\,\mathrm{d}C}
\newcommand{\dD}{\,\mathrm{d}D}
\newcommand{\dE}{\,\mathrm{d}E}
\newcommand{\dF}{\,\mathrm{d}F}
\newcommand{\dG}{\,\mathrm{d}G}
\newcommand{\dH}{\,\mathrm{d}H}
\newcommand{\dI}{\,\mathrm{d}I}
\newcommand{\dJ}{\,\mathrm{d}J}
\newcommand{\dK}{\,\mathrm{d}K}
\newcommand{\dL}{\,\mathrm{d}L}
\newcommand{\dM}{\,\mathrm{d}M}
\newcommand{\dN}{\,\mathrm{d}N}
\newcommand{\dO}{\,\mathrm{d}O}
\newcommand{\dP}{\,\mathrm{d}P}
\newcommand{\dQ}{\,\mathrm{d}Q}
\newcommand{\dR}{\,\mathrm{d}R}
\newcommand{\dS}{\,\mathrm{d}S}
\newcommand{\dT}{\,\mathrm{d}T}
\newcommand{\dU}{\,\mathrm{d}U}
\newcommand{\dV}{\,\mathrm{d}V}
\newcommand{\dW}{\,\mathrm{d}W}
\newcommand{\dX}{\,\mathrm{d}X}
\newcommand{\dY}{\,\mathrm{d}Y}
\newcommand{\dZ}{\,\mathrm{d}Z}

\let\oldforall\forall
\let\oldexists\exists
\renewcommand{\forall}{\,\oldforall\,}
\renewcommand{\exists}{\,\oldexists\,}

% "Geschmackssache.", sagte der Affe und biss in die Seife.

\let\subob\leq
\let\supob\geq
\renewcommand{\leq}{\leqslant}
\renewcommand{\geq}{\geqslant}
\newcommand{\oldleq}{\subob}
\newcommand{\oldgeq}{\supob}

% \varepsilon is _way_ cooler than \epsilon.

\let\oldepsilon\epsilon
\renewcommand{\epsilon}{\varepsilon}
\let\eps\epsilon
\let\oldphi\phi
\renewcommand{\phi}{\varphi}
\let\oldtheta\theta
\renewcommand{\theta}{\vartheta}

% Either use \subseteq or \subsetneq, \subset is not needed and misguiding.
% bei Nick gibts kein \subseteq, das ist einfach notational overload

%\renewcommand{\subset}{\subseteq}
%\renewcommand{\supset}{\supseteq}
\newcommand{\sub}{\subseteq}

% correctly typesets := and =:
\newcommand{\defeq}{\mathrel{\vcentcolon=}}
\newcommand{\eqdef}{\mathrel{=\vcentcolon}}

% Creates a version of \mid that can (only) be used between
% \left and \right and scales correctly.
\newcommand{\mmid}{\,\middle|\,}

% Dots for parameters, to be used if f(\cdot) doesn't cut it
\newcommand{\pardot}{{}\mathbin{\cdot}{}}

\newcommand{\smallboxtimes}{\mathbin{\XBox}}

\renewcommand{\tilde}[1]{\widetilde{#1}}
\renewcommand{\hat}[1]{\widehat{#1}}


\let\longimplies\implies
\renewcommand{\implies}{\longimplies}

\newcommand{\myqedsymbol}{\ensuremath{\Box}}

% In case a proof doesn't work out as planned (requires tikz)

\newcommand{\semiqed}{
  \parfillskip=0pt\nobreak\hfill
  {
    \scalebox{0.5}{
    \def\sqedlength{0.25}
    \def\sqedlength{0.5}
    \begin{tikzpicture}[overlay]
      \draw[densely dashed](0.0, 0.0) -- (0.0, \sqedlength);
      \draw[densely dashed](0.0, \sqedlength) -- (\sqedlength, \sqedlength);
      \draw[densely dashed](\sqedlength, \sqedlength) -- (\sqedlength, 0.0);
      \draw[densely dashed](0.0, 0.0) -- (\sqedlength, 0.0);
    \end{tikzpicture}
      }
  }\hspace{2px}
  \par
}

\newcounter{proofstep}
\newcommand{\step}[1][]{
  \refstepcounter{proofstep}
  \ifthenelse{\equal{#1}{}}{%
    %omitted
    \par\textbf{\mystepname\,\theproofstep:}
  }{%supplied
    \par\textbf{\mystepname\,\theproofstep\,{\normalfont(#1)}:}\unskip
  }
}

%%% \vvvert (|||x||| Norm), copied from https://tex.stackexchange.com/a/193126
\DeclareFontFamily{U}{matha}{\hyphenchar\font45}
\DeclareFontShape{U}{matha}{m}{n}{
      <5> <6> <7> <8> <9> <10> gen * matha
      <10.95> matha10 <12> <14.4> <17.28> <20.74> <24.88> matha12
      }{}
\DeclareSymbolFont{matha}{U}{matha}{m}{n}
\DeclareFontSubstitution{U}{matha}{m}{n}

% Math symbol font mathb
\DeclareFontFamily{U}{mathx}{\hyphenchar\font45}
\DeclareFontShape{U}{mathx}{m}{n}{
      <5> <6> <7> <8> <9> <10>
      <10.95> <12> <14.4> <17.28> <20.74> <24.88>
      mathx10
      }{}
\DeclareSymbolFont{mathx}{U}{mathx}{m}{n}
\DeclareFontSubstitution{U}{mathx}{m}{n}

% Symbol definition
\DeclareMathDelimiter{\vvvert}{0}{matha}{"7E}{mathx}{"17}


%

\DeclareRobustCommand{\bbDelta}{% blackboard bold Delta
	\mathchoice{\Delta\llap{\scalebox{0.775}{$\Delta$\hspace{0.025em}}}}%
	{\Delta\llap{\scalebox{0.775}{$\Delta$\hspace{0.025em}}}}%
	{\Delta\llap{\scalebox{0.775}{$\scriptstyle\Delta$\hspace{0.025em}}}}%
	{\Delta\llap{\scalebox{0.775}{$\scriptscriptstyle\Delta$\hspace{0.025em}}}}
	}

\DeclareRobustCommand{\bbGamma}{% blackboard bold Gamma
	\mathchoice{\raisebox{\depth}{\scalebox{1}[-1]{$\mathbb L$}}}%
	{\raisebox{\depth}{\scalebox{1}[-1]{$\mathbb L$}}}%
	{\raisebox{\depth}{\scalebox{1}[-1]{$\scriptstyle\mathbb L$}}}%
	{\raisebox{\depth}{\scalebox{1}[-1]{$\scriptscriptstyle\mathbb L$}}}
	}

% Symbol für die Cube Kategorie

\makeatletter
\newcommand{\bbsquaretikz}{%
    \mathrel{%
        \scalerel*{%
            \tikz[scale=0.23, line cap=round, baseline=-0.04ex]{%
                \draw[line width=.8pt, rounded corners=.1pt] (0,0) rectangle (1,1);
                \draw[line width=.3pt] (0.2,0) -- (0.2,1);
            }
        }{\square}
    }
}
\makeatother

\DeclareMathOperator*{\bbox}{\bbsquaretikz}

% Symbol für die lineare Cube Kategorie

\makeatletter
\newcommand{\bblsquaretikz}{%
    \mathrel{%
        \scalerel*{%
            \tikz[scale=0.23, line cap=round, baseline=-0.04ex]{%
                \draw[line width=.8pt, rounded corners=.1pt] (0,0) rectangle (1,1);
                \draw[line width=.3pt] (0.2,0.2) -- (0.2,1);
                \draw[line width=.3pt] (0.2,0.2) -- (1,0.2);
            }
        }{\square}
    }
}
\makeatother

\DeclareMathOperator*{\bblox}{\bblsquaretikz}

% Symbol für die Ecke

\makeatletter
\newcommand{\bbcornertikz}{%
    \mathrel{%
        \scalerel*{%
            \tikz[scale=0.23, line cap=round, baseline=-0.04ex]{%
                \draw[line width=.8pt] (0,0) -- (1,0);
                \draw[line width=.8pt] (0,0) -- (0,1);
                \draw[line width=.3pt] (0.2,0.2) -- (0.2,1);
                \draw[line width=.3pt] (0.2,0.2) -- (1,0.2);
            }
        }{\square}
    }
}
\makeatother

\DeclareMathOperator*{\bbcor}{\bbcornertikz}

%

%prisma eventuell so
\makeatletter
\newcommand{\bprism}{%
  \mathord{%
    \vcenter{\hbox{%
      \tikz[scale=0.25, line cap=round, line join=round]{%
        % back triangle
        \draw[line width=.35pt]
          (0.45,0.25) -- (1.35,0.25) -- (0.90,1.05) -- cycle;
        % connecting edges
        \draw[line width=.35pt]
          (0,0) -- (0.45,0.25);
        \draw[line width=.35pt]
          (0.90,0) -- (1.35,0.25);
        \draw[line width=.35pt]
          (0.45,0.80) -- (0.90,1.05);
        % front triangle
        \draw[line width=.35pt, rounded corners=.1pt]
          (0,0) -- (0.90,0) -- (0.45,0.80) -- cycle;
      }%
    }}%
  }%
}
\makeatother

%

\newcommand*\lon{%    % das umgedrehte \colon
        \nobreak
        \mskip6mu plus1mu
        \mathpunct{}%
        \nonscript
        \mkern-\thinmuskip
        {:}%
        \mskip2mu
        \relax
}

%

% Compact labelled 3-cube macro.
%
% Usage:
% \smallcube{Name}
%   {top-left}{top-right}
%   {upper-left}{upper-right}
%   {middle-left}{middle-right}
%   {bottom-left}{bottom-right}
%
% The macro draws at the current TikZ origin, so place it inside a scope:
% \begin{scope}[shift={(2.5,0)}]
%   \smallcube{C_3^4}{12}{123}{12}{13}{02}{023}{012}{013}
% \end{scope}

\tikzset{
  smallcube edge/.style={
    draw=black!55,
    line width=0.35pt,
    line cap=round,
    shorten <=1.5pt,
    shorten >=1.5pt
  },
  smallcube label/.style={
    font=\scriptsize,
    inner sep=0.7pt,
    fill=white
  },
  smallcube name/.style={
    font=\scriptsize\bfseries,
    anchor=south
  }
}

\newcommand{\smallcube}[9]{%
  \coordinate (smallcubeBL) at (0,0);
  \coordinate (smallcubeBR) at (1.45,0);
  \coordinate (smallcubeML) at (0.48,0.68);
  \coordinate (smallcubeMR) at (1.93,0.68);
  \coordinate (smallcubeUL) at (0,1.36);
  \coordinate (smallcubeUR) at (1.45,1.36);
  \coordinate (smallcubeTL) at (0.48,2.04);
  \coordinate (smallcubeTR) at (1.93,2.04);

  \node[smallcube name] at (0.96,2.35) {$#1$};

  \draw[smallcube edge] (smallcubeTL) -- (smallcubeTR);
  \draw[smallcube edge] (smallcubeTL) -- (smallcubeML);
  \draw[smallcube edge] (smallcubeTR) -- (smallcubeMR);
  \draw[smallcube edge] (smallcubeUL) -- (smallcubeTL);
  \draw[smallcube edge] (smallcubeUL) -- (smallcubeUR);
  \draw[smallcube edge] (smallcubeUL) -- (smallcubeBL);
  \draw[smallcube edge] (smallcubeUR) -- (smallcubeTR);
  \draw[smallcube edge] (smallcubeBL) -- (smallcubeML);
  \draw[smallcube edge] (smallcubeBL) -- (smallcubeBR);
  \draw[smallcube edge] (smallcubeBR) -- (smallcubeUR);
  \draw[smallcube edge] (smallcubeBR) -- (smallcubeMR);
  \draw[smallcube edge] (smallcubeML) -- (smallcubeMR);

  \node[smallcube label] at (smallcubeTL) {$#2$};
  \node[smallcube label] at (smallcubeTR) {$#3$};
  \node[smallcube label] at (smallcubeUL) {$#4$};
  \node[smallcube label] at (smallcubeUR) {$#5$};
  \node[smallcube label] at (smallcubeML) {$#6$};
  \node[smallcube label] at (smallcubeMR) {$#7$};
  \node[smallcube label] at (smallcubeBL) {$#8$};
  \node[smallcube label] at (smallcubeBR) {$#9$};
}

%

\allowdisplaybreaks[2] % encourages pagebreaks in formulas multiple lines long
                       % NOTE: Does not work with align, alignat, ... but e.g.
                       % with align*.

\theoremstyle{plain}

\newtheorem{theorem}{Theorem}[subsection]
\newtheorem{lemma}[theorem]{Lemma}
\newtheorem{corollary}[theorem]{Corollary}
\newtheorem{proposition}[theorem]{Proposition}
\newtheorem{conjecture}[theorem]{Conjecture}

% German stuff

\newtheorem{satz}[theorem]{Satz}
\newtheorem{folgerung}[theorem]{Folgerung}

\theoremstyle{definition}

\newtheorem{reminder}[theorem]{Reminder}
\newtheorem{remark}[theorem]{Remark}
\newtheorem{warning*}[theorem]{Warning}
\newtheorem{notation}[theorem]{Notation}
\newtheorem{note}[theorem]{Note}
\newtheorem{motivation}[theorem]{Motivation}
\newtheorem{assumption}[theorem]{Assumption}

\theoremstyle{definition}

\newtheorem{definition}[theorem]{Definition}
\newtheorem{question}[theorem]{Question}
\newtheorem{problem}[theorem]{Problem}

\newtheorem{exercise}{Exercise}[subsection]

\newtheorem{example}[theorem]{Example}
\newtheorem{construction}[theorem]{Construction}

% unnumbered environments

\theoremstyle{plain}

\newtheorem*{theorem*}{Theorem}
\newtheorem*{lemma*}{Lemma}
\newtheorem*{corollary*}{Corollary}
\newtheorem*{proposition*}{Proposition}
\newtheorem*{conjecture*}{Conjecture}

\newtheorem*{satz*}{Satz}
\newtheorem*{folgerung*}{Folgerung}

\theoremstyle{remark}

\newtheorem*{reminder*}{Reminder}
\newtheorem*{remark*}{Remark}
\newtheorem*{notation*}{Notation}
\newtheorem*{motivation*}{Motivation}
\newtheorem*{assumption*}{Assumption}

\theoremstyle{definition}

\newtheorem*{definition*}{Definition}
\newtheorem*{question*}{Question}
\newtheorem*{problem*}{Problem}
\newtheorem*{exercise*}{Exercise}
\newtheorem*{example*}{Example}
\newtheorem*{construction*}{Construction}

\usetikzlibrary{calc}
\usetikzlibrary{arrows.meta}

\newcommand{\bigrtimes}{%
  \mathop{\scalerel*{\rtimes}{\sum}}\displaylimits
}

%%
%% Theorems etc. in der Einleitung nach hinten referenzieren
%%

\makeatletter

\newenvironment{repeatresult}[2]{%
  \begingroup

  % Save environment name and label for the closing part
  \def\repeatresult@env{#1}%
  \def\repeatresult@label{#2}%

  % Save theorem machinery
  \let\repeatresult@origbegintheorem\@begintheorem
  \let\repeatresult@origopargbegintheorem\@opargbegintheorem
  \let\repeatresult@origrefstepcounter\refstepcounter

  % Do not increment a counter for the repeated statement
  \def\refstepcounter##1{}%

  % hyperref patches theorem counters separately
  \@ifundefined{Hy@theorem@refstepcounter}{}{%
    \let\repeatresult@orighyrefstepcounter\Hy@theorem@refstepcounter
    \def\Hy@theorem@refstepcounter##1{}%
  }%

  % Replace the displayed theorem number by the referenced number
  \def\@begintheorem##1##2{%
    \repeatresult@origbegintheorem
      {##1}{\ref{\repeatresult@label}}%
  }%

  \def\@opargbegintheorem##1##2##3{%
    \repeatresult@origopargbegintheorem
      {##1}{\ref{\repeatresult@label}}{##3}%
  }%

  % Start theorem / lemma / proposition / ...
  \csname #1\endcsname

  % Restore normal machinery for the theorem body
  \let\refstepcounter\repeatresult@origrefstepcounter

  \@ifundefined{repeatresult@orighyrefstepcounter}{}{%
    \let\Hy@theorem@refstepcounter\repeatresult@orighyrefstepcounter
  }%

  \let\@begintheorem\repeatresult@origbegintheorem
  \let\@opargbegintheorem\repeatresult@origopargbegintheorem
}{%
  \csname end\repeatresult@env\endcsname
  \endgroup
}

\makeatother

